%%%
%
% $Autor: Siddhanth Sharma $
% $Datum: 2021-05-14 $
% $Dateiname: 
% $Version: 4620 $
%
% !TeX spellcheck = GB
% !TeX program = pdflatex
% !TeX encoding = utf8
%
%%%

\chapter{\textcolor{red}{Domain Knowledge -- Tools}}

\index{Domain Knowledge -- Tools}


\section{Development Environment -- Python}

\subsection{Version}
\noindent
The project was developed using the Python programming language as the primary software tool on the NVIDIA Jetson Nano B01 (4~GB).  
Python is included by default in the JetPack SDK provided by NVIDIA and was therefore used as the system Python interpreter. The exact Python version was verified directly on the Jetson Nano during the initial setup phase using the command \texttt{python3 --version}. This ensured compatibility with the JetPack software stack and the installed computer vision and machine learning libraries.

\subsection{Description}
\noindent
Python forms the core development environment for the license plate detection system. It is used throughout the entire processing pipeline, including camera interfacing, image acquisition, preprocessing, model inference, and result handling. The choice of Python was motivated by its extensive ecosystem for computer vision and machine learning, as well as its seamless integration with the NVIDIA Jetson platform.

\vspace{1mm}
\noindent
The following Python libraries are central to the implementation:
\begin{itemize}
	\item \textbf{OpenCV (cv2)} for camera access, image preprocessing, and visualization tasks.
	\item \textbf{Ultralytics YOLOv8} for real-time license plate detection using a deep learning--based object detection model.
	\item \textbf{NumPy} for efficient numerical computations and array operations.
	\item \textbf{SQLite (sqlite3)} for lightweight, file-based data storage and logging of detection results.
\end{itemize}

\noindent
Together, these libraries provide a compact and efficient Python-based development environment suitable for real-time inference on an embedded edge device.

\subsection{Installation}
\noindent
Python was available on the Jetson Nano as part of the JetPack SDK installation and did not require a separate manual installation.  
All required Python packages were installed using the \texttt{pip} package manager associated with the system Python interpreter.

\vspace{1mm}
\noindent
The installation process was performed directly on the Jetson Nano to ensure binary compatibility with the ARM64 architecture. Required dependencies for computer vision, numerical processing, and model inference were installed sequentially using \texttt{pip}. This approach ensures a reproducible and lightweight software environment tailored to the target hardware.

\subsection{Configuration}
\noindent
No additional Python-specific configuration was required beyond the default JetPack setup.  
The CUDA, cuDNN, and TensorRT components provided by JetPack were automatically detected by the installed Python libraries where applicable. All Python scripts were executed using the system environment without modifying default library paths or environment variables.

\subsection{First Steps}
\noindent
After completing the installation, the Python environment was validated through a series of initial verification steps. These steps included importing core libraries such as OpenCV and NumPy and executing short test scripts to confirm correct installation and runtime behavior. Successful execution without import or runtime errors indicated that the Python environment was correctly configured and ready for further development.


\noindent
Figure~\ref{fig:python_verification_flow} summarizes the verification workflow used to confirm that the Python environment and the required core libraries are operational on the Jetson Nano before executing the main project pipeline.

\begin{figure}[h]
	\centering
	\begin{tikzpicture}[
		node distance=10mm,
		box/.style={draw, rounded corners, align=center, inner sep=6pt, text width=0.78\linewidth},
		arrow/.style={-Latex, thick}
		]
		\node[box] (jetson) {Jetson Nano (JetPack installed)};
		\node[box, below=of jetson] (python) {System Python Interpreter (\texttt{python3})};
		\node[box, below=of python] (deps) {Installed Dependencies via \texttt{pip}\\(e.g., OpenCV, NumPy, Ultralytics YOLOv8, SQLite)};
		\node[box, below=of deps] (hello) {Python Verification Script\\(\texttt{opencv\_version\_check.py})};
		\node[box, below=of hello] (ok) {Successful Execution\\(library import without errors, version printed)};
		
		\draw[arrow] (jetson) -- (python);
		\draw[arrow] (python) -- (deps);
		\draw[arrow] (deps) -- (hello);
		\draw[arrow] (hello) -- (ok);
	\end{tikzpicture}
	\caption{Python environment verification flow on the Jetson Nano.}
	\label{fig:python_verification_flow}
\end{figure}

\subsection{Program ``Hello World''}
\noindent
As an initial functionality test of the Python development environment, a minimal verification program was executed.  
The program imports the OpenCV library and prints the installed OpenCV version to the console. Successful execution confirms that Python and the core computer vision library required for the project are correctly installed and operational on the Jetson Nano.

\lstinputlisting[
language=Python,
caption={Python ``Hello World'' program for verifying OpenCV installation},
label={lst:python_hello_world_opencv},
basicstyle=\ttfamily\small,
keywordstyle=\color{blue},
commentstyle=\color{gray},
stringstyle=\color{green!40!black},
breaklines=true,
showstringspaces=false,
numbers=left,
numberstyle=\tiny\color{gray},
frame=single,
rulecolor=\color{black!30}
]{../../Code/OpenCV/opencv_version_check.py}



% -------------------------------------------------------------------------
% SECTION 2: Development Environment -- Hardware / OS / System Tools
% -------------------------------------------------------------------------
% NOTE:
% This section belongs under:
% \part{Domain Knowledge -- Tools}
% and comes after Section 1 (Python).
% -------------------------------------------------------------------------

\section{Development Environment -- Hardware / OS / System Tools}

\subsection{Version}
\noindent
The license plate detection system was deployed on an NVIDIA Jetson Nano B01 (4~GB) as the target edge device.  
The operating system and NVIDIA software stack were installed via the JetPack SDK (Linux for Tegra, L4T).  
All relevant system versions were verified directly on the Jetson Nano to ensure reproducibility and compatibility with GPU-accelerated inference.

\vspace{1mm}
\noindent
The following version information should be documented as part of the final system fingerprint:
\begin{itemize}
	\item \textbf{Hardware:} Jetson Nano B01 (4~GB).
	\item \textbf{JetPack SDK:} JetPack version installed on the device (verified via \SHELL{dpkg} or JetPack release information).
	\item \textbf{OS / L4T:} Ubuntu-based L4T release (verified using \SHELL{/etc/nv\_tegra\_release} and \SHELL{lsb\_release -a}).
	\item \textbf{CUDA Toolkit:} CUDA version (verified using \SHELL{nvcc --version}).
	\item \textbf{TensorRT:} TensorRT version provided by JetPack (verified using \SHELL{dpkg -l | grep tensorrt}).
\end{itemize}

\subsection{Description}
\noindent
This project uses the Jetson Nano as an embedded GPU-enabled edge platform to execute real-time computer vision workloads locally.  
The system tools and operating environment support the following core requirements:
\begin{itemize}
	\item \textbf{Stable OS base:} Ubuntu (L4T) as provided through JetPack for hardware-specific drivers and CUDA support.
	\item \textbf{Edge inference capability:} NVIDIA GPU acceleration through CUDA and TensorRT for efficient on-device inference.
	\item \textbf{Developer workflow:} development and deployment via terminal and remote access (e.g., SSH), enabling iterative testing and reproducible execution.
\end{itemize}

\noindent
In the project workflow, system-level tools are used primarily to (i) install and maintain the JetPack stack, (ii) validate hardware availability (CPU/GPU/memory), (iii) monitor runtime status, and (iv) execute verification and deployment scripts.


\noindent
Figure~\ref{fig:system_access_workflow} illustrates the typical access and verification workflow for the Jetson Nano system environment.

\begin{figure}[h]
	\centering
	\begin{tikzpicture}[
		node distance=10mm,
		box/.style={draw, rounded corners, align=center, inner sep=6pt, text width=0.78\linewidth},
		arrow/.style={-Latex, thick}
		]
		\node[box] (flash) {Flash JetPack/L4T Image to microSD Card};
		\node[box, below=of flash] (boot) {First Boot Initialization\\(user setup, network, updates)};
		\node[box, below=of boot] (access) {Device Access\\GUI (monitor/keyboard) \quad or \quad Headless (SSH)};
		\node[box, below=of access] (verify) {System Verification\\OS/L4T, CUDA, TensorRT, \texttt{tegrastats}};
		\node[box, below=of verify] (dev) {Run Project Pipeline\\(Python scripts for camera + inference + logging)};
		
		\draw[arrow] (flash) -- (boot);
		\draw[arrow] (boot) -- (access);
		\draw[arrow] (access) -- (verify);
		\draw[arrow] (verify) -- (dev);
	\end{tikzpicture}
	\caption{System setup and access workflow for the Jetson Nano development environment.}
	\label{fig:system_access_workflow}
\end{figure}

\subsection{Installation}
\noindent
The base system installation was performed by flashing a JetPack-compatible image onto a microSD card and booting the Jetson Nano from this card.  
JetPack provides the operating system (L4T), device drivers, and the NVIDIA acceleration libraries required for embedded computer vision workloads.

\vspace{1mm}
\noindent
A reproducible installation procedure typically includes:
\begin{itemize}
	\item \textbf{SD card flashing:} write the JetPack/L4T image to a microSD card using a flashing tool (e.g., balenaEtcher).
	\item \textbf{First boot initialization:} initial user creation, locale/time settings, and network setup (Ethernet/Wi-Fi).
	\item \textbf{System update:} update package lists and apply security updates via \SHELL{apt}.
	\item \textbf{Toolchain verification:} confirm that CUDA/TensorRT components provided by JetPack are present and accessible.
\end{itemize}

\noindent
This installation approach ensures that the Jetson Nano uses NVIDIA-supported drivers and libraries, which is essential for stable GPU-accelerated inference and camera interfacing.

\subsection{Configuration}
\noindent
After installation, the system was configured to support reliable development and execution of the project pipeline.  
Only configurations relevant to the operation of an embedded inference device were applied.

\vspace{1mm}
\noindent
Typical configuration steps include:
\begin{itemize}
	\item \textbf{Remote access (recommended):} enable and verify SSH access for development from a separate workstation.
	\item \textbf{System resources:} verify available storage and memory; configure swap only if required for model conversion or dependency compilation.
	\item \textbf{Performance monitoring:} ensure system monitoring tools (e.g., \SHELL{tegrastats}) are available for runtime observation.
	\item \textbf{Camera support:} verify that the selected camera interface (USB or CSI) is recognized by the OS before running the detection pipeline.
\end{itemize}

\noindent
All configuration changes should be kept minimal and documented, since the Jetson Nano is a resource-constrained embedded platform and relies on a stable baseline JetPack installation.

\subsection{First Steps (GUI / UI)}
\noindent
After installation and configuration, the first interaction with the device consists of validating the operating environment and confirming basic device functionality.  
Depending on the chosen workflow, the Jetson Nano can be used either with a directly connected monitor/keyboard (GUI) or remotely (headless) via SSH.

\vspace{1mm}
\noindent
Recommended first steps for verification are:
\begin{itemize}
	\item \textbf{Login and terminal access:} open a shell either locally (GUI terminal) or remotely (SSH).
	\item \textbf{System identity check:} confirm OS and L4T release information.
	\item \textbf{Acceleration stack check:} verify CUDA and TensorRT availability.
	\item \textbf{Runtime monitoring:} start \SHELL{tegrastats} to confirm real-time reporting of CPU/GPU load, memory usage, and temperature.
	\item \textbf{Camera presence (if applicable):} confirm that the camera device is detected by the OS before executing project scripts.
\end{itemize}

\subsection{Program ``Hello World'' / Bash Verification}
\noindent
To validate that the Jetson Nano system environment is operational, a minimal OS-level verification script was executed.  
The script prints core system information and performs basic checks that are directly relevant for edge inference workloads (OS identity, CUDA availability, TensorRT presence, and system monitoring).


 \lstinputlisting[
 language=bash,
 caption={OS-level ``Hello World'' verification script for Jetson Nano},
 label={lst:jetson_hello_world_bash},
 basicstyle=\ttfamily\small,
 keywordstyle=\color{blue},
 commentstyle=\color{gray},
 stringstyle=\color{green!40!black},
 breaklines=true,
 showstringspaces=false,
 numbers=left,
 numberstyle=\tiny\color{gray},
 frame=single,
 rulecolor=\color{black!30}
 ]{../../Code/SystemTools/jetson_hello_world.sh}